%

%\documentclass[11pt,a4paper,twoside]{article}

\usepackage[utf8]{inputenc}

\usepackage[T1,T2A]{fontenc}

\usepackage[english.ancient,russian.ancient]{babel}

\usepackage{graphicx}

\usepackage{array}

\usepackage[doi]{samgtu-bib}

\usepackage{samgtu}

\usepackage{samgtu-name}

\usepackage{mathtools}

\mathtoolsset{showonlyrefs}

\usepackage{desclist}

\usepackage{euscript}

\usepackage{mathrsfs} 

\usepackage{multicol}

\usepackage{mathrsfs}

\usepackage{cite}

\usepackage{cmap}

\usepackage{qrcode}

\allowdisplaybreaks[4]

\usepackage[nodayofweek]{datetime}

\usepackage[thinlines]{easybmat}



\renewcommand{\JournalYear}{2018}

\renewcommand{\JournalVolume}{22}

\renewcommand{\JournalNumber}{4}



\newdate{JournalDateSign}{29}{12}{2018}% Дата подписания Вестника в печать

\renewcommand{\JournalDateSign}{\displaydate{JournalDateSign}}

\newdate{JournalDatePrint}{16}{01}{2019}% Дата выхода Вестника из печати

\renewcommand{\JournalDatePrint}{\displaydate{JournalDatePrint}}

%\renewcommand{\JournalUPL}{16.56} % Усл. печ. л.

%\renewcommand{\JournalUIL}{16.53} % Уч.-изд. л.

\renewcommand{\JournalUPL}{15.24} % Усл. печ. л.

\renewcommand{\JournalUIL}{15.21} % Уч.-изд. л.

\renewcommand{\JournalRegNumber}{220/18} % Регистрационный номер

\renewcommand{\JournalOrderNumber}{2} % Номер заказа



%\renewcommand{\JournalRMonth}{Март} % Месяц выхода Вестника

%\renewcommand{\JournalEMonth}{March} % Месяц выхода Вестника

%\renewcommand{\JournalRMonth}{Июнь} % Месяц выхода Вестника

%\renewcommand{\JournalEMonth}{June} % Месяц выхода Вестника

%\renewcommand{\JournalRMonth}{Сентябрь} % Месяц выхода Вестника

%\renewcommand{\JournalEMonth}{September} % Месяц выхода Вестника

\renewcommand{\JournalRMonth}{Декабрь} % Месяц выхода Вестника

\renewcommand{\JournalEMonth}{December} % Месяц выхода Вестника



\newcommand{\totop}[1]{\raisebox{6pt}[1pt][2pt]{#1}}

\newcommand{\tottop}[1]{\raisebox{12pt}[1pt][2pt]{#1}} 

\newcommand{\totttop}[1]{\raisebox{18pt}[1pt][2pt]{#1}} 

\newcommand{\tottttop}[1]{\raisebox{24pt}[1pt][2pt]{#1}} 

\newcommand{\totttttop}[1]{\raisebox{30pt}[1pt][2pt]{#1}} 

\newcommand{\tobot}[1]{\raisebox{-6pt}[1pt][2pt]{#1}} 

\newcommand{\tobbot}[1]{\raisebox{-12pt}[1pt][2pt]{#1}} 

\newcommand{\tobbbot}[1]{\raisebox{-18pt}[1pt][2pt]{#1}}



\graphicspath{{./00PIC/}}

%

%\usepackage{samgtu-name}

%\usepackage{wrapfig}

%

%%\samgtupubl % для генерации pdf, отдаваемого в типографию СамГТУ

%\pdfcrop % обрезка pdf для размещения на math-net.ru

%\draft % На корректуре

%

%

\vsgtu{1517}

%

\FirstPage{122}

\LastPage{136}

%

%\begin{document}

%

%

\ResearchArticle

%%\ReviewArticle

%%\ShortCommunication

%

%

\renewcommand{\baselinestretch}{0.9}

%

%\hfill \qrcode[hyperlink,height=.8in]{http://doi.org/10.14498/vsgtu1517}

%\vspace{-.8in}

%





\udc{517.958:539.3}

\msc{74F15, 74B20, 74Q99}



\rutitle{Мезоскопические модели для определения упругих характеристик 

мягких магнитных эластомеров}

\rutitleshort{Мезоскопические модели для определения упругих характеристик\dots}



\ruauthor{А.~М.~Биллер, О.~В.~Столбов}{Б\,и\,л\,л\,е\,р~А.~М.,\ С\,т\,о\,л\,б\,о\,в~О.~В.}



\ruaffil{\rucmmperm.}

\enaffil{\encmmperm.}





\ruauthorsdetails{2}{

\fullauthor{\rupid{122212}{Анастасия Михайловна Биллер}} 

\CAuthor

\orcid{0000-0003-0355-4749}

\position[n]{младший научный сотрудник} 

\dept{лаб. физики и механики мягкого вещества}

\email[br]{kam@icmm.ru}

\fullauthor{\rupid{125086}{Олег Валерьевич Столбов}} 

\orcid{0000-0001-9088-7909}

\regalia{кандидат физико-математических наук} 

\position{научный сотрудник} 

\dept{лаб. физики и~механики мягкого вещества}

\email{sov@icmm.ru}}



\enauthorsdetails{2}{

\fullauthor{\enpid{122212}{Anastasiya M. Biller}} 

\CAuthor

\orcid{0000-0003-0355-4749}

\position[n]{Juniour Researcher} 

\dept{Physics and Mechanics of Soft Matter Laboratory}

\email{kam@icmm.ru}

\fullauthor{\enpid{125086}{Oleg V. Stolbov}} 

\orcid{0000-0001-9088-7909}

\regalia{Cand. Phys. \& Math. Sci.} 

\position{Researcher} 

\dept{Physics and Mechanics of Soft Matter Laboratory}

\email[b]{sov@icmm.ru}}





\rukeywords{мягкий магнитный эластомер, магнитомеханика, магнитомеханический гистерезис, магнитоиндуцированное упрочнение}



\ruabstract{В качестве модели мезоскопического структурного элемента мягкого магнитного эластомера рассмотрена пара намагничивающихся частиц, заключенных в цилиндр из высокоэластичного (гиперупругого) материала. 

Под действием внешнего магнитного поля частицы намагничиваются, и между ними возникает силовое взаимодействие. Частицы меняют свое положение в эластомерной матрице, насколько это позволяет ее упругое сопротивление. Равновесное положение частиц в образце определяется балансом магнитных и упругих сил и соответствует минимуму общей энергии системы. 

При ее вычислении учитывались как нелинейность и неоднородность намагничивания частиц, так и нелинейность упругих свойств эластомера. Это существенно приближает нас к~случаю реального магнитореологического композита, представляющего собой мягкий эластомер наполненный микронными частицами ферромагнетика. Рассматриваемая система проявляет бистабильность~--- при увеличении и уменьшении приложенного магнитного поля расстояние между частицами сокращается гистерезисным образом, от величины порядка нескольких радиусов до плотного контакта (коллапс). Такое поведение значительно влияет на способность мезоскопического элемента сопротивляться внешней нагрузке. 

Коллапс частиц под действием магнитного поля или сжимающей нагрузки приводит к~резкому росту жесткости. 

Зависимость механических характеристик от приложенного магнитного поля изучена для элементов различной относительной податливости. Эта зависимость также имеет гистерезисный характер. Несмотря на свою простоту, модель в~целом верно описывает индуцированную полем перестройку внутренней структуры мягких магнитных эластомеров. Полученные результаты использованы для качественного рассмотрения макроскопической магнитомеханики композита, это сделано с помощью процедуры гомогенизации, основанной на гипотезе Фойгта. Полученная зависимость жесткости мягкого магнитного эластомера от внешнего магнитного поля качественно совпадает с экспериментальными результатами, опубликованными в~литературе.}



\entitle{Mesoscopic models for definition of the large-scale elastic properties of the soft magnetic elastomers}

\entitleshort{Mesoscopic models for definition of the large-scale elastic properties\dots}



\enauthor{A.~M.~Biller, O.~V.~Stolbov}{B\,i\,l\,l\,e\,r~A.~M., S\,t\,o\,l\,b\,o\,v~O.~V.}



\enabstract{A pair of magnetizing particles embedded in a cylinder made of a high-elasticity (hyperelastic) material is considered as a model of a mesoscopic structure element of a soft magnetic elastomer. In the presence of the magnetic field particles magnetize and the force interaction is arisen between them. Particles change position inside the elastomer matrix as the elastic resistance allows it. Equilibrium position of the particles inside the sample is determined by the balance of magnetic and elastic forces and corresponds to the minimum of total energy of the system. In its calculation both the non-linearity and heterogeneity of the magnetization of the particles and non-linearity of the elastic properties of the elastomer have been taken into account. This brings us to the real magnetorheological composite, that is a soft elastomer filled with a micron ferromagnetic particles. The considered system exhibits bistability: increase and decrease of the applied magnetic field, leads to change of the distance between the particles in hysteretic manner, from a few radii to the tight contact (collapse). This behavior significantly affects the ability of a mesoscopic element to resist external load. Collapse of the particles inside it by a magnetic field or compressive load causes sharp increase of stiffness. The dependence of mechanical characteristics of the system on the strength of an applied magnetic field is studied for the elements of different compliance. This dependence also has a hysteresis. Despite its simplicity, the model in a generally correct way describes the field-induced changes of the internal structure of soft magnetic elastomers. The obtained results are used for qualitative analysis of the macroscopic magnetomechanics of the composite, this is done with the aid of a homogenisation procedure based of Voigt's hypothesis. The obtained dependence of the magnetic stiffness of soft magnetic elastomer on the external magnetic field agrees qualitatively with the published experimental results.}



\enkeywords{soft magnetic elastomer, magnetomechanics, magnetomechanical hysteresis, magnetoinduced stiffening}





\newdate{billerda}{25}{10}{2016}

\date{\displaydate{billerda}}

\newdate{billerdb}{10}{03}{2017}

\revisiondate{\displaydate{billerdb}}

\newdate{billerdc}{13}{03}{2017}

\accepteddate{\displaydate{billerdc}}

\newdate{billerdo}{03}{04}{2017}

\renewcommand{\JournalDataOnline}{\displaydate{billerdo}}







%%%%%%%%%%%%%%%%%%%%%%%%%%%%%%%%%%%%%%%%%%%%%%%%%%%%%%%%%%%%%%%%%%%%%%%%%%%%%%%%%%%%





\makerutitle





\Section[n]{Введение}

Мягкие магнитные эластомеры (ММЭ) выделяются среди других композитных материалов своей способностью к сильному изменению как формы, так и упругих свойств в ответ на приложение магнитного поля. 

Первопричиной возникновения всех этих особенностей являются микрочастицы магнитомягкого ферромагнетика, которыми наполнена высокоэластичная матрица ММЭ. Силовое взаимодействие частиц, вызванное их намагничиванием, заставляет ММЭ проявлять интересные эффекты~--- магнитоиндуцированные деформацию и упрочнение, эффект магнитной памяти формы, которые имеют большой практический потенциал. На сугубо качественном уровне картина проста. 

Приложенное внешнее поле намагничивает частицы и между их индуцированными магнитными моментами возникают пондермоторные силы, которые стремятся организовать частицы в структуру, отвечающую минимуму магнитостатической энергии. Однако в ММЭ этой тенденции противодействуют упругие силы (высокоэластичность), возрастающие по мере удаления частиц от их исходных положений внутри сетки. Чем выше степень сшивки полимерной матрицы, тем выше ее модуль упругости, а значит, тем сильнее она может ограничивать смещения частиц. Однако, как показывают теоретические исследования [\citen{Stolbov_11}--\citen{Han_IJSS_13}], даже небольшие изменения в пространственном распределении частиц внутри эластомера могут приводить к качественно различному поведению и свойствам моделируемого образца композита. Таким образом, критически важным становится понимание магнитоупругого взаимодействия частиц в мягком магнитном эластомере на мезоструктурном уровне как основы эффектов, проявляемых ММЭ.



\smallskip



%\section{\label{sec:2} Мезоскопический элемент}%%%	%%%



\Section{Мезоскопический элемент}

Ввиду важности межчастичных взаимодействий для моделирования и предсказания свойств ММЭ необходимо располагать как можно более точным представлением о магнитных и упругих силах, возникающих в композите на мезоскопическом уровне. Для изучения   этих вопросов в наших работах [\citen{JAP_14}--\citen{VMSS_15}] использовалась модельная система, представляющая собой пару сферических намагничивающихся частиц, помещенных в эластомерную матрицу (см.~рис.~\ref{fig:ME}). 

Такой мезоскопический образец является наименьшим представительным элементом ММЭ, позволяющим изучать его магнитомеханику.



\pagebreak



\begin{wrapfigure}{O}{.4\textwidth}

\centering

\includegraphics[scale=.9]{Element} 

\caption{Схематическое изображение мезоскопического элемента

[Figure~\ref{fig:ME}. The mesoscopic element schematic representation]

\label{fig:ME}}



%\vspace{-1mm}



\end{wrapfigure}



\vspace{-4mm}





В настоящей работе мы модифицировали нашу модель следующим образом. Размер эластомерного цилиндра, внутрь которого заключены две нелинейно намагничивающихся частицы, подобран так, что частицы занимают 30\% геометрического объема образца модельного ММЭ. Именно такова объемная концентрация частиц наполнителя в композитах, представляющих наибольший интерес для приложений. 

Частицы размещены вдоль оси цилиндра; они имеют сферическую форму и~одинаковый радиус $a$. Внешнее магнитное поле с~напряженностью ${\boldsymbol H}_0$ приложено вдоль оси образца. 

Под влиянием поля частицы намагничиваются и вступают в магнитостатическое (пондеромоторное) взаимодействие.

Магнитная энергия пары $U_{\text{mag}}$ в такой конфигурации зависит от расстояния между их центрами $l=|{\boldsymbol l}|$ и величины приложенного поля $H_0=|{\boldsymbol H}_0|$.  

Поскольку энергия пары достигает минимума при плотном контакте намагничивающихся частиц, действующие между ними силы "--- это притяжение, стремящееся передвинуть частицы навстречу друг другу. 

В~то же время смещение частиц вызывает деформацию эластомерной матрицы, в которую они встроены на условиях полной адгезии. 

Накапливаемая в матрице упругая энергия $U_{\text{el}}$ зависит от начального $l_0$ и текущего $l$ расстояний между центрами частиц; она минимальна в отсутствие деформаций матрицы, то есть при $l=l_0$. 

Таким образом, под действием внешнего магнитного поля в представительном элементе эластомера возникают два типа сил, конкурирующих между собой: магнитные, стремящиеся сблизить частицы, и упругие, противодействующие любой деформации матрицы.



Общая энергия $U$ рассматриваемого цилиндра с двумя частицами, включающая магнитную $ U_{\text{mag}}$ и упругую $U_{\text{el}}$ составляющие, имеет вид

\begin{equation}

\label{eq:U}%Общая энергия общая запись

U(l,l_0,H_0) =  U_{\text{mag}}(l,H_0) + U_{\text{el}}(l,l_0).

\end{equation}

В целом алгоритм получения общей энергии мезоскопического элемента в настоящей работе аналогичен тому, что описан в работе \cite{VMSS_15}. Приведем его коротко, сосредоточив внимание на выполненной модификации модели и на характерных свойствах рассматриваемой системы.



Для магнитной части $U_{\text{mag}}$ общей энергии представительного элемента использовались данные численного решения задачи магнитостатики для двух сферических частиц, которое было получено в работе \cite{PRE_15}. В этом решении учитывались как неоднородность, так и нелинейность намагничивания частиц. Связь между намагничиванием $\boldsymbol M$ и напряженностью поля внутри сферы $\boldsymbol H$ определялась эмпирическим законом Фрелиха--Кеннелли \cite{Bozorth_93}:

\begin{gather}

\label{eq:F-K}%Закон Фрелиха-Кеннелли

M({\boldsymbol H})=\frac{\chi_0 M_s {\boldsymbol H}}{M_s+\chi_0  H } 

=\left\{

\begin{array}{ll}

\chi_0 {\boldsymbol H} \quad & \text{при}\;\;\chi_0 H< M_s,\\

M_s \left(1 - M_s/\chi_0 H \right) \quad & \text{при}\;\; \chi_0 H> M_s.

\end{array} \right.

\end{gather}

Здесь $\chi_0$ "--- начальная магнитная восприимчивость частиц, а $M_s$ "--- намагниченность насыщения. 

Для частиц карбонильного железа, чаще всего используемых в~качестве наполнителя для ММЭ, характерные значения этих величин составляют $10^4$ и $1500$~кА/м соответственно. 

В дальнейшем изложении мы будем использовать безразмерные единицы: 

напряженность магнитного поля ${h_0=H_0/M_s}$, 

расстояние между центрами частиц $q=l/a$ и магнитную энергию $\tilde{U}_{\text{mag}}(q,h_0)=U_{\text{mag}}(l,H_0)/(\mu_0 M_s^2 a^3)$.



Упругая составляющая $U_{\text{el}}$ энергии представительного элемента, накапливаемая в эластомерной матрице при перемещении частиц, была получена аналогично тому, как это было сделано в работе \cite{VMSS_15}, но расчеты выполнялись для цилиндрического образца эластомера с указанными выше пропорциями. 

Задача нелинейной теории упругости формулировалась в осесимметричной постановке и решалась методом конечных элементов. Поскольку эластомер может подвергаться значительным деформациям, его определяющее соотношение было выбрано в форме закона Муни"--~Ривлина для несжимаемой среды \cite{Oswald_09}:

\begin{equation}

\label{eq:M-R}%Закон Муни-Ривлина

W(\mathbf C)= c_1 \tilde W(\mathbf C)= c_1\left[ \left( I_1(\mathbf C) - 3\right) + \tilde c_2\left( I_2(\mathbf C) - 3 \right) \right],

\end{equation}

где $I_1(\mathbf C)$ и $I_2(\mathbf C)$ "--- инварианты меры деформации Коши"--~Грина $\mathbf C$, $c_1$ и $c_2$  "--- числовые константы модели. 

Здесь также введен обезразмеренный потенциал $\tilde W = W/c_1$ и параметр $\tilde c_2=c_2/c_1$.



В расчетах расстояние $l$ между частицами изменялось дискретно от начального $l_0$ до $l=2a$, соответствующего плотному контакту частиц. 

На каждом шаге методом конечных элементов определялось поле перемещений, затем вычислялась объемная плотность энергии. 

В~результате был сформирован дискретный массив данных, представляющий зависимость упругой энергии $U_{\text{el}}$ от начального $l_0$ и~текущего $l$ расстояний между центрами частиц.



Для построения и верификации интерполяционной формулы, описывающей зависимость $U_{\text{el}}$ от $l$, использовалась эвристическая схема, родственная той, что была предложена в работе \cite{VMSS_15}. 

Наша конструкция представляет собой систему цилиндрических стержней, которая по общей длине и по расстоянию между частицами идентична пропорциям расчетного образца. 

Нумерация стержней в этой схеме начинается с {\sl 2}; номер {\sl 1} зарезервирован для обозначения в дальнейшем относительного изменения расстояния между частицами.



\begin{figure}[b!]

\centering \footnotesize

\includegraphics[scale=.9]{st_model} 

\caption{Эвристическая схема для интерполяции упругой энергии

\label{fig:Scheme}}



[Figure~\ref{fig:Scheme}. The heuristic scheme for the elastic energy interpolation]



\end{figure}



Стержни подчиняются тому же соотношению Муни"--~Ривлина \eqref{eq:M-R}, что и эластомерная матрица в~численном расчете. 

Коэффициентами интерполяции служат радиусы стержней $r_k$ и длина $l_0^{(2)}$ внутреннего стержня {\sl 2}. 

Формула для $U_{\text{el}}$ при каждом конкретном $q_0$ имеет вид

\begin{equation}

\label{eq:Uel(q)}%

U_{\text{el}}(l)  = \pi c_1 \sum_{k=2}^4 (2-\delta_{k4}) \tilde W({\bf C}_k) r_k^2 l_0^{(k)}.

\end{equation}

Интерполяция строилась с учетом того, что при малых деформациях образца, вызванных перемещением частиц внутри него или внешней нагрузкой, разложение энергии в ряд совпадает с таким же разложением энергии стержневой модели. Этот факт позволяет найти часть интерполяционных коэффициентов по коэффициентам разложения, а оставшиеся подобрать из соображений наилучшего соответствия между энергией стержневой модели и расчетного образца в области больших деформаций.



На следующем шаге решения коэффициенты формулы \eqref{eq:Uel(q)} интерполируются посредством различных степенных функций. В результате мы получили непрерывную функцию

\begin{align}                                                          \label{eq:Uel}%

U_{\text{el}}(l,l_0) &= \pi c_1 \sum_{k=2}^4 (2-\delta_{k4}) \tilde W({\bf C}_k) r_k^2(l_0) l_0^{(k)}(l_0),\\

r_2 &= 1.003-0.2255/(1+e^{(0.8989 l_0^2-6.311)}),\nonumber\\

r_3 &= 1.103 - 0.01318 l_0^2 - 3.301 \cdot 10^{-6} l_0^8,\nonumber\\

r_4 &= 0.1637+0.02213 (l_0-3.045)^2,\nonumber\\

l_0^{(2)} &= 1.464 - 4.232 l_0 + 2.314 l_0^2 - 0.2781 l_0^3,\nonumber\\

l_0^{(3)} &= (l_0^{(4)}- l_0-2a)/2, \qquad l_0^{(4)} = 6.4a,\nonumber

\end{align}

описывающую упругую энергию, накапливаемую в эластомерном цилиндре с~погрешностью не более 5\%.



\begin{figure}[b!]

\centering \footnotesize

\includegraphics[scale=.9]{Hyst_1_9} 

\caption{Магнитомеханический гистерезис в системе с $\beta=110$ и начальным межцентровым расстоянием $q_0$, равным 2.5  (линия {\sl 1}\/), 3 (линия {\sl 2\/}), 3.5 (линия {\sl 3\/}), 4  (линия {\sl 4}\/) 

\label{fig:Hyster_9}}



\smallskip



[Figure 3. Magnetomechanical hysteresis in the system with $\beta=110$ and initial interparticle  distance $q_0$ equals 2.5  (line {\sl 1}\/), 3 (line {\sl 2\/}), 3.5 (line {\sl 3\/}), and 4  (line {\sl 4}\/)]



\end{figure}



Суммирование интерполяционных формул для магнитной и упругой составляющих энергии дает аналитическую функцию

\begin{equation}

\label{eq:U}%Общая энергия общая запись

U(q,q_0,h_0)/(c_1 a^3) =  \beta \tilde U_{\text{mag}}(q,h_0) + \tilde U_{\text{el}}(q,q_0),

\end{equation}

являющуюся хорошим приближением точной магнитоупругой энергии мезоскопического элемента; в ней относительная величина магнитного и упругого вкладов определяется параметром $\beta=\mu_0 M_s^2/c_1$.



Анализ функции \eqref{eq:Uel} показывает, что в некотором диапазоне полей система переходит в режим бистабильности: в ней одновременно существуют два положения устойчивого равновесия. Следствием бистабильности является возникновение гистерезиса межчастичного расстояния при изменении величины приложенного магнитного поля (см.\ рис.\ \ref{fig:Hyster_9}). В системах, где частицы намагничиваются с насыщением, возможность появления гистерезиса при заданном начальном расстоянии между частицами $q_0=l_0/a$ контролируется параметром $\beta$, который имеет смысл эффективной податливости композита. Отметим, что гистерезисное поведение возникает вне зависимости от того, находятся ли частицы в неограниченной эластомерной среде \cite{PRE_15} или в образце конечного размера \cite{VMSS_15}.







%\section{\label{sec:3} Мезоскопический элемент под действием внешних сил} %%%	%%%



\smallskip



\Section{Мезоскопический элемент под действием внешних сил} 

В случае отсутствия сторонних механических сил интерполяционная формула \eqref{eq:Uel} для энергии $U_{\text{el}}$ цилиндрического образца позволяет найти зависимость изменения его общего (<<внешнего>>) размера $u_4$ от изменения параметра $u_1$ "--- расстояния между центрами частиц по оси цилиндра.



В присутствии сторонней механической нагрузки (растяжение или сжатие, приложенное к торцам цилиндра) относительные изменения расстояния между частицами $\lambda_1$ и общей длины образца $\lambda_4$ становятся независимыми переменными. 

Соответственно, энергетическая функция  превращается в  $U(\lambda_1,\lambda_4,q_0,h_0)$, что дает возможность получить упругие характеристики модельного образца эластомера.



Введем систему координат так, чтобы ось $Oz$ была направлена вдоль вектора $\boldsymbol l$, соединяющего центры частиц в образце, упругий отклик которого мы описываем системой стержней (см. рис.~\ref{fig:Scheme}). 

Под действием внешнего магнитного поля, направленного по $Oz$, образец деформируется так, что относительные изменения межчастичного расстояния (воображаемый стержень {\sl 1}) и стержня {\sl 4} принимают значения $\varepsilon_1^*=\lambda_1^*-1$ и $\varepsilon_4^*=\lambda_4^*-1$ соответственно. 

В~отсутствие внешней нагрузки состояние системы, задаваемое параметрами $(\varepsilon_1^*,\varepsilon_4^*)$, является равновесным.



В зависимости от величины приложенного поля равновесному состоянию системы могут соответствовать как небольшое сближение частиц, так и их фактический коллапс: $q\simeq2$. Рассмотрим, к примеру, мезоскопический элемент с $q_0=3.5$, магнитомеханический гистерезис которого в числе других представлен на рис.~\ref{fig:Hyster_9}. Такой гистерезис является <<латентным>> \cite{PRE_15} в том смысле, что, хотя парный кластер является одним из двух возможных устойчивых положений системы, перейти в это состояние только под действием магнитного поля она не может. Однако коллапс частиц можно реализовать при добавлении внешней механической нагрузки, и это, конечно, отразится на общей длине образца. Приложим к торцам образца силы величиной $f$, сжимающие или растягивающие его. 

Изменяя величину приложенной силы и вычисляя равновесную относительную деформацию мезоскопического элемента $\varepsilon_4$, то есть решая задачу оптимизации $U(\varepsilon_1,\varepsilon_4) - fu_4 \to \text{min}$, 

получаем диаграмму <<напряжение--деформация>>, которая представлена на рис.~\ref{fig:T(E)_9},~{\sl a}.



\begin{figure}[h!]

\centering \footnotesize



\includegraphics[scale=.9]{T(E)_9} 



\caption{Диаграммы <<напряжение--деформация>> мезоскопических элементов с относительной податливостью $\beta=110$ и~начальными расстояниями между центрами частиц:

 {\sl a})~${q_0=3.5}$ в~поле с~напряженностью $h_0=0.32$;  

 {\sl b}) $q_0=3.25$ в полях c~напряженностью $h_0=0.16$ (линия {\sl 1}\/), 0.24 (линия {\sl 2}\/), 0.32 (линия {\sl 3}\/)

\label{fig:T(E)_9}}



\smallskip



[Figure~\ref{fig:T(E)_9}. Stress--strain diagrams for mesoscopic elements with relative compliance $\beta=110$ and initial interparticle distance: {\sl a\/}) $q_0=3.5$ under applied field intensity $h_0=0.32$ and {\sl b}\/)~$q_0=3.25$ under applied field intensity $h_0=0.16$ (line {\sl 1}\/), 0.24 (line {\sl 2}\/), and 0.32 (line {\sl 3}\/)]



\vspace{-3mm}

\end{figure}



Как видно из рис.~\ref{fig:T(E)_9},~{\sl a}, если мезоскопический элемент намагничен, то он уже немного деформирован даже в отсутствие внешней механической силы (см.~точку {\sl 1}). 

Это сжатие соответствует известному магнитострикционному эффекту, присущему ММЭ. 

Пусть к~образцу приложена сжимающая нагрузка, стремящаяся уменьшить его длину. 

С~ростом силы $f$ сжатие образца будет сначала идти плавно, но по достижении силой некоторого конечного значения произойдет коллапс частиц, и деформация изменится скачком (см.~точку {\sl 2}\/). 

На диаграмме рис.~\ref{fig:T(E)_9},~{\sl a} это соответствует перескоку представляющей точки с~нижней ветки кривой на верхнюю и дальнейшему движению вдоль последней. 

Примечательно, что при снятии нагрузки система не возвращается в начальное положение (см.~точку {\sl 3}\/), а~продолжает оставаться в~сжатом состоянии (<<магнитная скрепка>>). 

Разрушить эту структуру можно только с помощью усилия противоположного направления "--- растягивающего (см.~точку~{\sl 4}\/). 

В~результате графическое представление цикла <<нагрузка--разгрузка>> отображается на рис.~\ref{fig:T(E)_9},~{\sl a} в~виде петли, то есть на определенном участке функция обладает двузначностью.



Приведенные рассуждения показывают, что если в мезоскопическом элементе существует магнитомеханический гистерезис, то на кривой <<нап\-ря\-же\-ние--деформация>> петля будет появляться в магнитном поле любой напряженности. 

Этот вывод иллюстрирует серия кривых на рис.~\ref{fig:T(E)_9},~{\sl b}, представляющая деформационные циклы для мезоскопического элемента с~$q_0=3.25$ и~$\beta=110$. 

Как видно, положение петли зависит от величины приложенного поля. 

Так, для элемента, находящегося в поле с~напряженностью $h_0=0.16$, она целиком расположена в~области сжимающей нагрузки, а~для того же элемента в~поле $h_0=0.32$  петля имеет место только при растягивающих усилиях. Из этого следует, что, регулируя величину приложенного магнитного поля, можно моделировать свойства мезоскопического элемента: устанавливать величину сил, при которых будет происходить коллапс или расцепление частиц, то есть увеличиваться или уменьшаться способность элемента сопротивляться деформации.



Необходимо отметить, что эффективная податливость мезоскопического элемента $\beta=110$ (ед.\ СИ) соответствует достаточно жесткой полимерной матрице. Для реальных ММЭ с модулем Юнга порядка 10--30~кПа 

(см.~например \cite{Stepanov_07}) величина параметра $\beta$ достигает нескольких тысяч. 

В~этом случае в~рассматриваемых системах с~начальным межчастичным расстоянием $q_0$ от 2 до 4 магнитомеханический гистерезис существует всегда, и~коллапс частиц происходит в полях меньших $h_0=0.1$ (см.~рис.~\ref{fig:Hyster_200}).



\begin{figure}[h!]

\centering \footnotesize

\includegraphics[scale=.9]{Hyst_1_200} 

\caption{Магнитомеханический гистерезис в~системе с~$\beta=2500$ и начальным межцентровым расстоянием $q_0$, равным 2.5 (линия~{\sl 1\/}), 3 (линия~{\sl 2\/}), 3.5 (линия~{\sl 3\/}), 4 (линия~{\sl 4\/}) 

\label{fig:Hyster_200}}



\smallskip



[Figure~\ref{fig:Hyster_200}. Magnetomechanical hysteresis in the system with $\beta=2500$ and initial interparticle  distance $q_0$ equals  2.5 (line {\sl 1}\/), 3 (line {\sl 2}\/), 3.5 (line {\sl 3}\/), and 4 (line {\sl 4}\/)]







\end{figure}



Разрушить <<магнитную скрепку>>, образовавшуюся в элементе с высокой податливостью, непросто. 

Для этого требуется растянуть образец в~несколько раз. Для примера в таблице приведены характерные значения параметра $f/(c_1S)$, необходимые для образования и разрушения <<магнитной скрепки>> в~мезоскопическом элементе с~относительной податливостью $\beta=2500$, а~также значения деформации элемента до ($\varepsilon_s$) и после ($\varepsilon_e$) 

по обе стороны от скачкообразного перехода. 

Данные в таблице соответствуют точкам, отмеченным на рис.~\ref{fig:T(E)_9},~{\sl a}: {\sl 2} "--- коллапс частиц и {\sl 4} "--- разрушение парного кластера.



\begin{table}[tbh!]

\centering \small

{Характерные точки диаграммы <<напряжение--деформация>> для мезоскопического образца с $\beta=2500$ и~расстоянием между центрами частиц $q_0=3.5$ }



[Typical points of the stress--strain diagram for the mesoscopic element with $\beta=2500$ and interparticle distance $q_0=3.5$]



%\renewcommand{\arraystretch}{1.8}

\renewcommand{\tabcolsep}{0.28cm}

\begin{tabular}{c|c|c|c|c|c|c|c}

\hline    

		& \footnotesize $f/(c_1S)=0$	&\multicolumn{3}{c|}{ \footnotesize Point 2, see Fig.~\ref{fig:T(E)_9},~{\sl a}}				

		&\multicolumn{3}{c}{\footnotesize Point 4, see Fig.~\ref{fig:T(E)_9},~{\sl a}} 				\\ \hline                                       

\footnotesize $h_0$		& \footnotesize $\varepsilon$	&\footnotesize $f/(c_1S)$	&\footnotesize $\varepsilon_s$	&\footnotesize $\varepsilon_e$	&\footnotesize $f/(c_1S)$	& \footnotesize $\varepsilon_s$	& \footnotesize$\varepsilon_e$	\\ \hline

\rule{0mm}{12pt}%

$0.02$	& $-0.021$		& $-1.58$	& $-0.109$		& $-0.140$ 		&22.95	&\phantom{$1$}5.494 		&\phantom{$1$}3.981		\\ 

$0.04$	& $-0.080$		& \phantom{$-$}3.40		&\phantom{$-$}0.416		&\phantom{$-$}0.285		&60.90	&14.727 		&11.288		\\ 

$0.06$	& $-0.080$		& \phantom{$1$}22.95	&\phantom{$-$}5.373		&\phantom{$-$}3.981		&95.40	&21.946 		&17.933		\\ \hline

\end{tabular}

\end{table}



%\section{\label{sec:4} Упругие свойства мезоскопического элемента} %%%	%%%





\Section{Упругие свойства мезоскопического элемента}

Для определения механических характеристик мезоскопического элемента разложим его энергию в~ряд в окрестности положения равновесия $(\varepsilon_1^*,\,\varepsilon_4^*)$:

\begin{gather}

\label{eq:U_series_1}

U(\varepsilon_1^*,\varepsilon_4^*)=\frac{1}{2}

\frac{\partial^2 U}{\partial \varepsilon_1^2}\Bigl|_{\scriptsize\begin{matrix}\scriptsize \varepsilon_4=\varepsilon_4^*, \\ \scriptsize\varepsilon_1=\varepsilon^*_1\end{matrix}}\varepsilon_1^2 +

\frac{1}{2}

\frac{\partial^2 U}{\partial \varepsilon_4^2}\Bigl|_{\scriptsize\begin{matrix} \varepsilon_4=\varepsilon_4^*, \\ \varepsilon_1=\varepsilon^*_1\end{matrix}}\varepsilon_4^2 +

\frac{\partial^2 U}{\partial \varepsilon_1 \partial \varepsilon_4}\Bigl|_{\scriptsize\begin{matrix} \varepsilon_4=\varepsilon_4^*, \\ \varepsilon_1=\varepsilon^*_1\end{matrix}}\varepsilon_1 \varepsilon_4.

\end{gather}

По понятным причинам производные первого порядка здесь обращаются \linebreak в~нуль.



Приложенная малая нагрузка вызывает дополнительную малую деформацию образца. Возникающие при этом силы должны уравновешивать как магнитные силы взаимодействия частиц, так и внешнюю механическую нагрузку. Следовательно, энергия образца в сумме с работой этих сил должна стремиться к минимуму:

\begin{equation}

\label{eq:OP}%

U(\varepsilon_1,\varepsilon_4) - fu_4 \to \text{min}.

\end{equation}

Запишем работу внешних сил в виде $fu_4=f\varepsilon_4 l_0^{(4)}$, используя соотношение между деформацией и~перемещением 

$\varepsilon_4=u_4/l_0^{(4)}$. 

Тогда  задача оптимизации \eqref{eq:OP} запишется в виде

\begin{equation}

\label{eq:U_series_2}%

U - f\varepsilon_4 l_0^{(4)} = \frac{1}{2} A_{11} \varepsilon_1^2 + \frac{1}{2} A_{22} \varepsilon_4^2 +  A_{12} \varepsilon_1 \varepsilon_4 - f\varepsilon_4 l_0^{(4)} \to \min,

\end{equation}

где $A_{ij}$ "--- вторые частные производные энергии.

В этих обозначениях экстремум рассматриваемой функции описывается решениями системы уравнений

\begin{equation}

\label{eq:Diff}%

\begin{array}{l}

\displaystyle

\frac{\partial U}{\partial \varepsilon_1} = A_{11} \varepsilon_1 + A_{12} \varepsilon_4 = 0, \\[3mm]

\displaystyle

\frac{\partial U}{\partial \varepsilon_4} = A_{22} \varepsilon_4 + A_{12} \varepsilon_1 = f l_0^{(4)}.

\end{array}

\end{equation}



Введем обозначение $f/\varepsilon_4=K$, которое определяет способность рассматриваемой модели сопротивляться деформации. 

Из системы \eqref{eq:Diff} видно, что 

\begin{equation}

\label{eq:K}%

K=\frac{f}{\varepsilon_4} = \Bigl( A_{22} - \frac{A_{12}^2}{A_{11}} \Bigr) \frac{1}{l_0^{(4)}}.

\end{equation}

Определим относительную жесткость мезоскопического образца $\tilde K$ как отношение найденной $K$ системы в магнитном поле к начальной жесткости $K_0$ в~его отсутствие. 



Равновесное положение частиц $(\varepsilon_1^*,\,\varepsilon_4^*)$ в элементе изменяется в зависимости от приложенного магнитного поля гистерезисным образом. Аналогично изменяется и жесткость рассматриваемой системы. 

На рис.~\ref{fig:Hyster_9} было представлено изменение расстояния между частицами в мезоскопических элементах с~различными $q_0$ под действием внешнего магнитного поля. 

Соответствующее изменение относительной жесткости элементов $\tilde K$ приведено на рис.~\ref{fig:E}.



\begin{figure}[b!]

\centering \footnotesize



\includegraphics[scale=.9]{K_q0_200} 

\caption{Изменение относительной жесткости мезоскопических элементов с $\beta=2500$ в зависимости от приложенного магнитного поля при начальном расстоянии между частицами~$q_0$, равным 2.5 (линия~{\sl 1\/}), 3 (линия~{\sl 2\/}), 3.5 (линия~{\sl 3\/}), 4 (линия~{\sl 4\/})  

\label{fig:E}}



\smallskip



[Figure 6. Relative stiffness of mesoscopic elements with $\beta=2500$ depend on the apllied mag\-ne\-tic field for the initial interparticle distance $q_0$ equals 2.5 (line {\sl 1}\/), 3 (line {\sl 2}\/), 3.5 (line {\sl 3}\/), and 4 (line {\sl 4}\/)]



\end{figure}



С ростом внешнего поля частицы плавно сближаются, а модуль системы уменьшается. 

Затем при определенной величине поля в системе появляется вторая жесткость, что соответствует возникновению в системе второго положения равновесия. С дальнейшим увеличением поля происходит коллапс частиц, который влечет за собой резкий рост жесткости мезоскопического элемента. Чем больше начальное расстояние между частицами $q_0$, тем сильнее эффект упрочнения двухчастичного образца в магнитном поле. 



\smallskip



%\section{\label{sec:5} Процедура гомогенизации}%%%	%%%





\Section{Процедура гомогенизации}

Двухчастичная модель позволяет ра\-зобраться в~основных закономерностях перестройки внутренней структуры ММЭ и~следующим за ней изменением свойств композита. 

Однако для предсказания свойств макроскопических образцов необходимо перенести на них результаты мезоскопического моделирования. 

Для этого используются процедуры гомогенизации,  некоторые варианты которых описаны в~работах \cite{Castaneda_11,Menzel_14}. 

Поскольку наша модель учитывает лишь одну конфигурацию частиц в~эластомере во внешнем поле, хотя и вносящую серьезный вклад, мы ограничимся здесь простой процедурой осреднения.



Представим некоторый макроскопический образец ММЭ как систему наших мезоскопических элементов, каждый из которых содержит пару частиц с~различным начальным межцентровым расстоянием $q_0$. 

Предположим, что к образцу приложена растягивающая сила в~направлении вдоль длин цилиндров. 

Будем считать, что каждый мезоскопический элемент под действием этой силы будет растягиваться независимо. 

Таким образом, мы не учитываем никакого взаимодействия на границах между отдельными элементами. 

Для вычисления модуля упругости макроскопического образца или напряжений, возникающих при его растяжении/сжатии, требуется только усреднить соответствующие характеристики по всем мезоскопическим элементам. 

Подобные модели материалов называют статистическими, основанными на гипотезе Фойгта. 



Описываемый подход можно сравнить с построением феноменологической модели материала, где элементами схемы выступают наши мезоскопические образцы, соединенные параллельно. Поскольку число частиц в макроскопических образцах ММЭ значительно, то и число элементов в схеме необходимо устремить к бесконечности и усреднение проводить по непрерывной функции распределения расстояния между центрами частиц. 

Таким образом, жесткость макроскопического образца с заданным распределением расстояний между частицами вычисляется по формуле

\begin{equation}                                                          

\label{eq:Esum}%

\langle K(H_0) \rangle = {\int_{q_0^{\text{min}}}^{q_0^{\text{max}}} K(q_0,h_0)f(q_0)dq_0}

\bigg/{\int_{q_0^{\text{min}}}^{q_0^{\text{max}}} f(q_0)dq_0},

\end{equation}

где $f(q_0)$ "--- функция плотности распределения $q_0$. 

Для ММЭ логично предположить нормальное распределение расстояний между частицами наполнителя. 

Максимальное межчастичное расстояние, которое допускает  наш мезоскопический образец, 

составляет $4.25a$. 

Зададим распределение со средним значением $m=3a$ и среднеквадратичным отклонением $\sigma=0.2$. 

В~результате получим кривую зависимости жесткости ММЭ от величины приложенного поля $h_0$, которая представляет собой простое усреднение множества подобных кривых для мезоскопических элементов. 



\begin{figure}[p]

\centering \footnotesize

\includegraphics[scale=.9]{Esum} 

\caption{Относительная жесткость образца ММЭ, имеющего нормальное распределение меж\-частичных \mbox{расстояний с $m\,{=}\,3$ и~$\sigma\,{=}\,0.2$ в зависимости от внешнего магнитного поля~$h_0$}

\label{fig:g_E}}



[Figure~\ref{fig:g_E}. Relative stiffness of the SME sample, that has normal distribution with $m=3$ and $\sigma=0.2$, depend on the extrenal magnetic field $h_0$]





\includegraphics[scale=.9]{Esum_other} 

\caption{Зависимость относительной жесткости образцов ММЭ от внешнего магнитного поля $h_0$. 

Образцы имеют нормальное распределение межчастичных расстояний: 

{\sl a}\/) с~различным $m$ и $\sigma=0.2$; вид распределений представлен на части {\sl b} рисунка; 

{\sl c}\/) с~$m=3$ и~различным $\sigma$; вид распределений представлен на части {\sl d} рисунка\label{fig:g_E_other}}



\smallskip

 

[Figure~\ref{fig:g_E_other}. Dependence of relative stiffness of the SME samples on the external magnetic field~$h_0$. Samples have normal distribution of interparticle distances: {\sl a}\/) with different $m$ and~$\sigma=0.2$; distributions presented on the part {\sl b} of the plot; 

{\sl c}\/) with $m=3$ and different $\sigma$; distributions presented on the part {\sl d} of the plot]





\end{figure}



Жесткость макроскопического образца ММЭ сначала слегка уменьшается, а затем плавно растет с увеличением поля и в полях порядка $\sim 0.6$ в~безразмерных единицах достигает насыщения. 

В обратном направлении при снятии внешнего поля наблюдается гистерезис  жесткости. 

Такое поведение легко объяснить, опираясь на графики изменения модуля для мезоскопических элементов с различным $q_0$, представленные на рис.~\ref{fig:E}. При включении поля частицы сближаются и во всех элементах без исключения жесткость уменьшается. 

По достижении напряженностью внешнего поля значения  $\simeq 0.02$ на графике рис.~\ref{fig:g_E} возникают два различных значения жесткости. 

Это соответствует области гистерезиса в мезоскопических элементах с~межчастичным расстоянием около~$2.5$. 

С~ростом поля первыми в~плотный контакт приходят пары с наименьшим $q_0$. 

Увеличение жесткости в~таких парах не настолько велико, чтобы уравновесить его уменьшение в~элементах с~большим $q_0$. 

В~рассматриваемом ММЭ наибольшее число пар расположено на расстоянии $q_0=3$ между центрами частиц, поэтому в~поле $h_0=0.045$, когда эти пары коллапсируют, начинается значительный рост жесткости образца. 

В~поле $h_0=0.06$ даже частицы, находящиеся на большом расстоянии друг от друга, схлопываются, рост жесткости ММЭ останавливается и эффект упрочнения достигает насыщения. 



Приведенное выше описание позволяет без труда разобраться, как вид распределения $q_0$ влияет на получаемую в результате жесткость макроскопического образца ММЭ. 

На рис.~\ref{fig:g_E_other} представлены зависимости $\tilde K$ от величины приложенного поля для других распределений. 

Смещение пика распределения в~область больших межчастичных расстояний приводит к увеличению максимальной жесткости ММЭ. 

При этом уменьшение разброса $q_0$ увеличивает разницу между наибольшим и наименьшим значением жесткости образца.





\smallskip

%\section{\label{sec:6} Заключение}%%%	%%%



\Section[n]{Заключение}

В настоящей работе показано, как гистерезисное поведение частиц в мезоскопическом образце влияет на его механические характеристики. На качественном уровне эффект магнитоиндуцированного упрочнения ММЭ может быть объяснен образованием плотных кластеров из частиц внутри композита при его намагничивании, что проиллюстрировано в работе с помощью процедуры осреднения. Несмотря на то, что рассматриваемая модель учитывает только те пары частиц, которые расположены вдоль приложенного магнитного поля, представленные зависимости относительной жесткости ММЭ от внешнего магнитного поля качественно совпадают с кривыми изменения упругих модулей структурированных образцов композита, полученных экспериментально и опубликованных в работах \cite{Stepanov_07,Abramchuk_06}. 





\AdditionalContent{Конкурирующие интересы}{Мы не имеем конкурирующих интересов.} 

\AdditionalContent{Авторский вклад и ответственность}{Все  авторы  принимали  участие  в  разработке  концепции  \mbox{статьи}  и  в  написании  рукописи. Авторы несут  полную  ответственность  за  предоставление  окончательной рукописи в печать. Окончательная  версия рукописи была одобрена всеми авторами.} 

\AdditionalContent{Финансирование}{Работа выполнена при поддержке Комплексной программы фундаментальных исследований УрО РАН №~10 (проект № 15--10--1--18) и при поддержке РФФИ (проекты №~16--31--00503 и №~17--42--590504).}











%% Благодарности, гранты и прочее



\RusRef



\begin{thebibliography}{99}



\Bibitem{Stolbov_11}

\paper Modelling of magnetodipolar striction in soft magnetic elastomers

\by Stolbov~O.~V., Raikher~Yu.~L., Balasoiu~M.

\jour Soft Matter

\yr  2011

\issue 18

\vol 7

\pages 8484--8487

\crossref{10.1039/c1sm05714f}





\Bibitem{Ivaneyko_12}

\paper Effects of particle distribution on mechanical properties of magneto-sensitive elastomers in a homogeneous magnetic field

\by Ivaneyko~D., Toshchevikov~V.~P., Saphiannikova~M., Heinrich~G.

\jour Condensed Matter Physics

\yr 2012

\issue 3

\vol 15

\pages 33601:1--12

\crossref{10.5488/cmp.15.33601}

\arxiv \href{https://arxiv.org/abs/1210.1401}{1210.1401} [cond-mat.soft]





\Bibitem{Han_IJSS_13}

\paper Field-stiffening effect of magneto-rheological elastomers

\by Han~Y., Hong~W., Faidley~L.~E.

\jour International Journal of Solids and Structures

\yr 2013

\issue 14--15

\vol 50

\pages 2281--2288

\crossref{10.1016/j.ijsolstr.2013.03.030}



\Bibitem{JAP_14}

\paper  Modeling of particle interactions in magnetorheological elastomers

\by Biller~A.~M., Stolbov~O.~V., Raikher~Yu.~L.

\jour J. Appl. Phys.

\yr 2014

\issue 11

\vol 116

\papernumber 114904 

\totalpages 8 \nofrills

\crossref{10.1063/1.4895980}



\Bibitem{PRE_15}

\paper Mesoscopic magnetomechanical hysteresis in a magnetorheological elastomer

\by Biller~A.~M., Stolbov~O.~V., Raikher~Yu.~L.

\jour Phys. Rev.~E

\yr 2015

\issue 2

\vol 92

\papernumber 023202 

\totalpages 9 \nofrills

\crossref{10.1103/physreve.92.023202}



\RBibitem{VMSS_15}

\paper Бистабильное магнитомеханическое поведение ферромагнитных частиц в эластомерной матрице

\by Биллер~А.~М., Столбов~О.~В., Райхер~Ю.~Л.

\jour Вычисл. мех. сплош. сред

\yr 2015

\issue 3

\vol 8

\pages 273--288

\elib{http://elibrary.ru/item.asp?id=24388654}

\crossref{10.7242/1999-6691/2015.8.3.23}



\Bibitem{Bozorth_93}

\by Bozorth~R.~M.

\book Ferromagnetism

\publaddr  Piscataway, N.J.

\publ Wiley-IEEE Press

\yr 1993,

%\totalpages 

xvii+969~pp

\crossref{10.1109/9780470544624}



\Bibitem{Oswald_09}

\by Oswald~P.

\book Rheophysics: The Deformation and Flow of Matter

\publaddr New York

\publ Cambridge University Press

\yr 2009

\totalpages 624



\Bibitem{Stepanov_07}

\paper Effect of a homogeneous magnetic field on the viscoelastic behavior of magnetic elastomers

\by Stepanov~G.~V., Abramchuk~S.~S., Grishin~D.~A., Nikitin~L.~V., Kramarenko~E.~Yu., Khokhlov~A.~R.

\jour Polymer

\yr 2007

\issue 2

\vol 48

\pages 488--495

\crossref{10.1016/j.polymer.2006.11.044}

 





\Bibitem{Castaneda_11}

\paper Homogenization-based constitutive models for magnetorheological elastomers at finite strain

\by Ponte Casta\~neda~P., Galipeau~E.

\jour  J. Mech. Phys. Solids

\yr 2011

\issue 2

\vol 59

\pages 194--215

\crossref{10.1016/j.jmps.2010.11.004}





\Bibitem{Menzel_14}

\paper  Bridging from particle to macroscopic scales in uniaxial magnetic gels

\by Menzel~A.

\jour J. Chem. Phys.

\yr 2014

\issue 19 

\vol 141

\papernumber 194907 

\totalpages 13 \nofrills

\crossref{10.1063/1.4901275}





\Bibitem{Abramchuk_06}

\paper  Effect of a homogeneous magnetic field on the mechanical behavior of soft magnetic elastomers under compression

\by Abramchuk~S.~S., Grishin~D.~A., Kramarenko~E.~Yu., Stepanov~G.~V., Khokhlov~A.~R.

\jour Polym. Sci. Ser. A

\yr 2006

\issue 2

\vol 48

\pages 138--145

\crossref{10.1134/s0965545x06020064}



\end{thebibliography}







\newpage





\makeentitle



\vspace{-3mm}



\AdditionalContent{Competing interests}{We have no competing interests.} 

\AdditionalContent{Authors' contributions and responsibilities}{Each author has participated in the article concept development and in the manuscript writing.

The authors are absolutely responsible for submitting the final manuscript in print.

Each author has approved the final version of manuscript.} 

\AdditionalContent{Funding}{This work was supported by the UB RAS Program no.~10 (project no.~15--10--1--18), the Russian Foundation for Basic Research (projects' nos. 16--31--00503, and 17--42--590504).}





%Competing interests. I declare I have no competing interests.

%Competing interests. We have no competing interests.

%Competing interests. We declare we have no competing interests.

%Acknowledgements. The authors acknowledge the referees for their positive comments and suggestions.

%Authors' contributions and responsibilities

%

%Исследование не имело спонсорской поддержки.

%Автор несет полную ответственность за предоставление окончательной версии рукописи в печать. Окончательная версия рукописи была одобрена автором.

%%Автор не получал гонорар за статью.

%The author is absolutely responsible for submitting the final manuscript in print.

%The author has approved the final version of manuscript.



%}







\vspace{-3mm}





\begin{thebibliography}{99}



\Bibitem{Stolbov_11:en}

\paper Modelling of magnetodipolar striction in soft magnetic elastomers

\by Stolbov~O.~V., Raikher~Yu.~L., Balasoiu~M.

\jour Soft Matter

\yr  2011

\issue 18

\vol 7

\pages 8484--8487

\crossref{10.1039/c1sm05714f}





\Bibitem{Ivaneyko_12:en}

\paper Effects of particle distribution on mechanical properties of magneto-sensitive elastomers in a homogeneous magnetic field

\by Ivaneyko~D., Toshchevikov~V.~P., Saphiannikova~M., Heinrich~G.

\jour Condensed Matter Physics

\yr 2012

\issue 3

\vol 15

\pages 33601:1--12

\crossref{10.5488/cmp.15.33601}

\arxiv \href{https://arxiv.org/abs/1210.1401}{1210.1401} [cond-mat.soft]





\Bibitem{Han_IJSS_13:en}

\paper Field-stiffening effect of magneto-rheological elastomers

\by Han~Y., Hong~W., Faidley~L.~E.

\jour International Journal of Solids and Structures

\yr 2013

\issue 14--15

\vol 50

\pages 2281--2288

\crossref{10.1016/j.ijsolstr.2013.03.030}



\Bibitem{JAP_14:en}

\paper  Modeling of particle interactions in magnetorheological elastomers

\by Biller~A.~M., Stolbov~O.~V., Raikher~Yu.~L.

\jour J. Appl. Phys.

\yr 2014

\issue 11

\vol 116

\papernumber 114904 

\totalpages 8 \nofrills

\crossref{10.1063/1.4895980}



\Bibitem{PRE_15:en}

\paper Mesoscopic magnetomechanical hysteresis in a magnetorheological elastomer

\by Biller~A.~M., Stolbov~O.~V., Raikher~Yu.~L.

\jour Phys. Rev.~E

\yr 2015

\issue 2

\vol 92

\papernumber 023202 

\totalpages 9 \nofrills

\crossref{10.1103/physreve.92.023202}



\Bibitem{VMSS_15:en}

\lang In Russian

\paper Bistable magnetomechanical behavior of ferromagnetic particles in an elastomer matrix 

\by Biller~A.~M., Stolbov~O.~V., Raikher~Yu.~L.

\jour Vychisl. mekh. splosh. sred \rm [Computational continuum mechanics]

\yr 2015

\issue 3

\vol 8

\pages 273--288

\crossref{10.7242/1999-6691/2015.8.3.23}



\Bibitem{Bozorth_93:en}

\by Bozorth~R.~M.

\book Ferromagnetism

\publaddr  Piscataway, N.J.

\publ Wiley-IEEE Press

\yr 1993,

%\totalpages 

xvii+969~pp

\crossref{10.1109/9780470544624}



\Bibitem{Oswald_09:em}

\by Oswald~P.

\book Rheophysics: The Deformation and Flow of Matter

\publaddr New York

\publ Cambridge University Press

\yr 2009

\totalpages 624



\Bibitem{Stepanov_07:en}

\paper Effect of a homogeneous magnetic field on the viscoelastic behavior of magnetic elastomers

\by Stepanov~G.~V., Abramchuk~S.~S., Grishin~D.~A., Nikitin~L.~V., Kramarenko~E.~Yu., Khokhlov~A.~R.

\jour Polymer

\yr 2007

\issue 2

\vol 48

\pages 488--495

\crossref{10.1016/j.polymer.2006.11.044}

 





\Bibitem{Castaneda_11:en}

\paper Homogenization-based constitutive models for magnetorheological elastomers at finite strain

\by Ponte Casta\~neda~P., Galipeau~E.

\jour  J. Mech. Phys. Solids

\yr 2011

\issue 2

\vol 59

\pages 194--215

\crossref{10.1016/j.jmps.2010.11.004}





\Bibitem{Menzel_14:en}

\paper  Bridging from particle to macroscopic scales in uniaxial magnetic gels

\by Menzel~A.

\jour J. Chem. Phys.

\yr 2014

\issue 19 

\vol 141

\papernumber 194907 

\totalpages 13 \nofrills

\crossref{10.1063/1.4901275}





\Bibitem{Abramchuk_06:en}

\paper  Effect of a homogeneous magnetic field on the mechanical behavior of soft magnetic elastomers under compression

\by Abramchuk~S.~S., Grishin~D.~A., Kramarenko~E.~Yu., Stepanov~G.~V., Khokhlov~A.~R.

\jour Polym. Sci. Ser. A

\yr 2006

\issue 2

\vol 48

\pages 138--145

\crossref{10.1134/s0965545x06020064}



\end{thebibliography}



\newpage



%\end{document}

%

